% !TeX spellcheck = en_US
\documentclass[final,hyperref={pdfpagelabels=false}]{beamer}
\usepackage{grffile}
\mode<presentation>
\usetheme{I6pd2}
\usepackage[english]{babel}
\usepackage[utf8]{inputenc}
\usepackage{amsmath,amsthm, amssymb}
\usepackage{amsfonts,dsfont}
\usepackage{multirow}
\usepackage{graphicx}
\usepackage{adjustbox}
\usepackage[flushleft]{threeparttable}
\usepackage{array}
\usepackage{anyfontsize}
\usepackage[utf8]{inputenc}
%\usepackage{times}\usefonttheme{professionalfonts}  % obsolete
%\usefonttheme[onlymath]{serif}
\boldmath
\usepackage[orientation=portrait,size=a0,scale=1.3,debug]{beamerposter}
% change list indention level
% \setdefaultleftmargin{3em}{}{}{}{}{}


%\usepackage{snapshot} % will write a .dep file with all dependencies, allows for easy bundling
\newcolumntype{C}[1]{>{\centering\let\newline\\\arraybackslash\hspace{0pt}}m{#1}}

\usepackage{array,booktabs,tabularx}
\newcolumntype{Z}{>{\centering\arraybackslash}X} % centered tabularx columns
\newcommand{\pphantom}{\textcolor{ta3aluminium}} % phantom introduces a vertical space in p formatted table columns??!!

\listfiles

%Imagem de fundo
% \usebackgroundtemplate{\centering
%         \begin{picture}(1000,3320)
% 			\includegraphics[width=40em]{images/arvore_itv.pdf}
%          \end{picture}
% }
        
%Transparencia dos blocos
%\addtobeamertemplate{block begin}{\pgfsetfillopacity{0.7}}{\pgfsetfillopacity{1}}

%%%%%%%%%%%%%%%%%%%%%%%%%%%%%%%%%%%%%%%%%%%%%%%%%%%%%%%%%%%%%%%%%%%%%%%%%%%%%%%%%%%%%%
\graphicspath{{figures/}}
 
%\title{HYPERSPECTRAL IMAGE INTERPRETATION BASED ON PARTIAL LEAST SQUARES}
\title{{Lorem ipsum dolor sit amet,
consectetur adipiscing elit. Vestibulum sed
nisi lobortis ex posuere venenatis.}}	
\author{Gabriel Garcia\inst{1}, Autor 2\inst{1} \inst{2}}
\institute[shortinst]{\inst{1} Federal University of Ouro Preto, Ouro Preto, Brazil       \\ %\and %
                      \inst{2} Vale Institute of Technology, Ouro Preto, Brazil           
}
\date[Friday, July 29, 2016]{Friday, July 29, 2016}

%%%%%%%%%%%%%%%%%%%%%%%%%%%%%%%%%%%%%%%%%%%%%%%%%%%%%%%%%%%%%%%%%%%%%%%%%%%%%%%%%%%%%%
\newlength{\columnheight}
\setlength{\columnheight}{89cm}


%%%%%%%%%%%%%%%%%%%%%%%%%%%%%%%%%%%%%%%%%%%%%%%%%%%%%%%%%%%%%%%%%%%%%%%%%%%%%%%%%%%%%%
 \begin{document}
	
\begin{frame}
  \begin{columns}
    % ---------------------------------------------------------%
    % Set up a column 
    \begin{column}{.5\textwidth}
      \begin{beamercolorbox}[center,wd=\textwidth]{postercolumn}
        \begin{minipage}[T]{.98\textwidth}  % tweaks the width, makes a new \textwidth
          \parbox[t][\columnheight]{\textwidth}{ % must be some better way to set the the height, width and textwidth simultaneously
            % Since all columns are the same length, it is all nice and tidy.  You have to get the height empirically
            % ---------------------------------------------------------%
            % fill each column with content 
            %\vspace{0.1cm}      
            
            \begin{block}{Abstract}
				\begin{itemize}
				\item The electrocardiogram (ECG) is the most widely used non-invasive technique in heart disease diagnoses
			\begin{itemize}
				\item The ECG is frequently used to detect cardiac rhythm abnormalities, otherwise known as arrhythmias
			\end{itemize}
        \item The aiming of this work is:
            
\end{itemize}
            \end{block}
%   
			\vfill
          	\begin{block}{Arrhythmia classification system}
	           	
	           
	           	\begin{figure}[!htb]
	           	
	           			\centering
	           			\includegraphics[width=15cm]{figures/diagram.png}
	           	
	           	\end{figure}
	           	
	           		           	
           	\end{block}
         
 			\vfill
          	\begin{block}{The Temporal Vectorcardiogram}
	           	
	           
	           	\begin{figure}[!htb]
	           		\begin{minipage}{0.50\linewidth}
	           			\centering
	           			\includegraphics[width=13cm]{figures/ecg.png}
	           		\end{minipage} \hfill
	           		\begin{minipage}{0.49\linewidth}
	           			\begin{itemize}
	           				\item This á é î system responsible for the pumping action of heartbeats is formed by a group of specialized cells that send electrical signals
	           				\item Electrocardiogram (ECG) is the most common process to record these electrical signals
	           			
	           				
	           			\end{itemize}

	           		\end{minipage}
	           	\end{figure}
	           	
	           		           	\begin{figure}[!htb]	           		
	           		           		\begin{minipage}{0.40\linewidth}
	           		           			\begin{itemize}
	           		           				\item We propose a new representation method based in the VCG
	           		           				\item Where we take two leads and we add the time as a new axis, defining the temporal VCG
	           		           				\item It has a more rich representation than the normal VCG from ECG signal.
	           		           			\end{itemize}
	           		           		\end{minipage}\hfill
	           		           		\begin{minipage}{0.59\linewidth}
	           		           			\centering
	           		           			\includegraphics[width=16cm]{figures/vcg3d.png}
	           		           		
	           		           		\end{minipage} 
	           		           	\end{figure}
           	\end{block} 
			\vfill
			
			%\vspace{0.45\textheight} 
			 
%
%            
          }
        \end{minipage}
      \end{beamercolorbox}
    \end{column}
    % ---------------------------------------------------------%
    % end the column

    % ---------------------------------------------------------%
    % Set up a column 
    %\vspace{1cm}
    \begin{column}{.5\textwidth}
      \begin{beamercolorbox}[center,wd=\textwidth]{postercolumn}
        \begin{minipage}[T]{.98\textwidth}
          \parbox[t][\columnheight]{\textwidth}{
            \begin{block}{Classification}
            	\begin{figure}[!htb]	           		
            		\begin{minipage}{0.45\linewidth}
						\begin{itemize}
							\item An accurate arrhythmia classification is desirable to correctly diagnose
							cardiac issues and in some cases, the early detection can save lives
							\item SVM has been widely used to classify arrhythmias
						
							\end{itemize}
            		\end{minipage}\hfill
            		\begin{minipage}{0.55\linewidth}
            			\centering
            			\includegraphics[scale=0.17]{figures/svm.png}	
            		\end{minipage} 
            	\end{figure}
            	
            	
            \end{block}
            \vspace{0.2cm}
\vfill
 \begin{block}{(MIT-BIH) Arrhythmia Database}
            	
						\begin{itemize}
							\item It is essential a representative database for a good training and fair evaluation 
							\item The Massachusetts Institute of Technology - Beth Israel Hospital Arrhythmia Database is one of the recommended databases by the Association for the Advancement of Medical Instrumentation
							
						\end{itemize}
            	
            	
            	
            \end{block}

\vfill          
            
            \vfill
            	\begin{block}{Results}
            		\begin{itemize}
            			\item Test results (DS2) of the best Complex Network parameters configuration, chosen by the most relevant (DS1) training for Temporal VCG.
            		\end{itemize}      		
            		
            		{\small
            			
            			\begin{table}[h!]
            				
            				\begin{center}
            					
            					\label{tab:results3d} 
            					\centering	
            					\begin{tabular}{|c|c|c|c|c|c|c|c|} 
            						\hline 
            						Filter & $m$ & $T_0$ & $T_Q$ & Acc & N Se/+P/FPR & S Se/+P/FPR & V Se/+P/FPR   \\ \hline \hline
            						%%
            						de Chazal's & 2.0 & 0.055  & 0.6 &90.2 &   94.0/  96.3/  29.6 &   28.2/  22.2/   3.8 &   84.8/  64.7/   3.1   \\ \hline			
            						%%
            						Common   & 8.0 & 0.005 & 0.9 &   91.0 &   96.1/  95.2/  39.8 &   13.8/  17.1/   2.5 &   75.6/  68.8/   2.3 \\ 			\hline
            						%%
            						\textbf{No Filter}& \textbf{2.0} & \textbf{0.055} & \textbf{0.6} &  \textbf{91.2} &   \textbf{95.0/  96.5/  27.9} &  \textbf{ 29.6/  26.4/   3.1} &   \textbf{85.1/  66.3/   3.0} \\             	\hline \hline 
            					\end{tabular}
            					\vspace{0.05cm}
            					\linebreak 
            					
%            					{\footnotesize {Test results of the best complex networks parameters configuration for Temporal VCG.
%            							
%            							 \textbf{Se}, \textbf{+P} and \textbf{FPR} stands for Sensitivity, Positive Predictive Value and False Positive Rate, respectively.}}	
            					
            				\end{center}
            			\end{table}}
            			\begin{itemize}
            				\item Test results (DS2) of the best Complex Network parameters configuration, chosen by the most relevant (DS1) training for VCG-2D.
            			\end{itemize}
            			{\small				
            				\begin{table}[h!]					
            					\begin{center}						
            						\label{tab:results2d} 
            						\centering	
            						\begin{tabular}{|c|c|c|c|c|c|c|c|} 
            							\hline 
            							Filter & $m$ & $T_0$ & $T_Q$ & Acc & N Se/+P/FPR & S Se/+P/FPR & V Se/+P/FPR  \\ \hline \hline
            							%
            							de Chazal's & 2.0 & 0.055 & 0.6 &   87.3 &   90.9/  95.8/  32.6 &   30.1/  20.0/   4.6 &   82.0/  53.2/   5.0 \\ 			\hline
            							%
            							\textbf{Common}   & \textbf{4.0} & \textbf{0.005} & \textbf{0.9} &   \textbf{90.8} &   \textbf{95.1/  95.7/  35.0} &   \textbf{29.3/  23.0/   3.8} &   \textbf{76.6/  73.6/   1.9} \\			\hline
            							%
            							No Filter& 2.0 & 0.055 & 0.6 &    89.2 &   93.1/  96.1/  30.3 &   25.4/  17.8/   4.5 &   82.2/  62.7/   3.4 \\ 		\hline \hline
            						\end{tabular} 
            						\vspace{0.15cm}
            						\linebreak 						
            						%{\footnotesize {Test results of the best complex networks parameters configuration for VCG-2D.}}							
            					\end{center}
            				\end{table}}
            \begin{itemize}
            	\item The temporal VCG have increased the global accuracy, and have better results classifying the N and S classes, when it is compared with the best result of the VCG-2D.
            	
            \end{itemize}
            
            
            \end{block}
            \vspace{0.2cm}
            
            	            
            \vfill		            
            \begin{block}{Conclusions}
            	
            	
              \begin{itemize}
                \item In this work we proposed a new 3D visualization of the vectorcardiogram, so called temporal vectorcardiogram, and complex network as feature extractor method with SVM based arrhythmia classification in ECG signals. 
               
                \item The method achieves greater accuracy classifying Ventricular Ectopic Beats, Supraventricular Ectopic Beats and Normal Beats.
                
                \item New techniques to extract 3D features from the Temporal VCG could be an interesting research direction.
                
                
              \end{itemize}
            \end{block}
           % \vspace{0.45\textheight} 
            %\vfill
          }
          % ---------------------------------------------------------%
          % end the column
        \end{minipage}
      \end{beamercolorbox}
    \end{column}
    % ---------------------------------------------------------%
    % end the column
  \end{columns}
  \vskip1ex
\end{frame}
\end{document}


%%%%%%%%%%%%%%%%%%%%%%%%%%%%%%%%%%%%%%%%%%%%%%%%%%%%%%%%%%%%%%%%%%%%%%%%%%%%%%%%%%%%%%%%%%%%%%%%%%%%
%%% Local Variables: 
%%% mode: latex
%%% TeX-PDF-mode: t
%%% End:
